\documentclass[11pt]{article}
\usepackage[utf8]{inputenc}
\usepackage[T1]{fontenc}
\usepackage{hyperref}
\usepackage{listings}
\usepackage{xcolor}
\usepackage{geometry}
\usepackage{color}

\geometry{a4paper, margin=1in}

\definecolor{codegreen}{rgb}{0,0.6,0}
\definecolor{codegray}{rgb}{0.5,0.5,0.5}
\definecolor{codepurple}{rgb}{0.58,0,0.82}
\definecolor{backcolour}{rgb}{0.95,0.95,0.92}

\lstdefinestyle{mystyle}{
    backgroundcolor=\color{backcolour},
    commentstyle=\color{codegreen},
    keywordstyle=\color{magenta},
    numberstyle=\tiny\color{codegray},
    stringstyle=\color{codepurple},
    basicstyle=\ttfamily\footnotesize,
    breakatwhitespace=false,
    breaklines=true,
    captionpos=b,
    keepspaces=true,
    numbers=left,
    numbersep=5pt,
    showspaces=false,
    showstringspaces=false,
    showtabs=false,
    tabsize=2
}

\lstset{style=mystyle}

\title{VVS Assignment 2: Comprehensive Testing of the Calendar Application}
\author{Rodrigo Neto}
\date{\today}

\begin{document}

    \maketitle

    \section{Introduction}
    This report describes the testing strategy and implementation for the \textit{Calendar} application. The goal was to provide a robust test suite covering multiple levels of the application: Unit Testing, Integration Testing (3rd party, REST, and Database), and End-to-End (E2E) testing, along with Continuous Integration setup.

    \section{Step-by-Step Solution}
    The following steps were taken to fulfill the assignment requirements:
    \begin{enumerate}
        \item \textbf{Environment Setup}: Added Maven dependencies for JUnit 5, Mockito, Selenium, WebDriverManager, and Spring Security Test.
        \item \textbf{SUT Adaptation}: Modified the source code to improve testability, specifically injecting \texttt{RestClient.Builder} in discovery providers and adding \texttt{setId} methods to model classes.
        \item \textbf{Unit Testing}: Developed tests for services using Mockito to isolate business logic from persistence.
        \item \textbf{3rd Party Integration}: Used \texttt{@RestClientTest} to mock external API responses and verify parsing logic.
        \item \textbf{REST API Integration}: Expanded test coverage to all controllers (\texttt{AuthController}, \texttt{MeetingController}, \texttt{DiscoveryController}, \texttt{CalendarController}, \texttt{ICalController}) using \texttt{MockMvc} and \texttt{@WebMvcTest}.
        \item \textbf{Security Testing}: Implemented dedicated security tests to verify authorization rules for anonymous and authenticated users.
        \item \textbf{Database Integration}: Verified the JPA layer using \texttt{@SpringBootTest} with an H2 database.
        \item \textbf{End-to-End Testing}: Created multiple Selenium journeys to verify full user flows, including registration, login, discovery, and meeting creation.
        \item \textbf{CI/CD Optimization}: Categorized all tests using JUnit 5 tags and configured a multi-tier GitHub Actions pipeline with manual triggers, concurrency controls, and artifact management.
    \end{enumerate}

    \section{Testing Details}

    \subsection{Unit Testing (Business Logic)}
    Targeting the \texttt{service} and \texttt{discover} packages.
    \begin{itemize}
        \item \textbf{UserServiceTest}: Tests user registration (success and duplicate username cases) and user retrieval.
        \item \textbf{MeetingServiceTest}: Tests proposing meetings, responding to invites, and copying events from discovery.
        \item \textbf{AppUserDetailsServiceTest}: Ensures Spring Security correctly loads user details.
        \item \textbf{ICalServiceTest}: Verifies iCal feed generation and special character escaping.
        \item \textbf{DiscoveryServiceTest}: Tests event aggregation, de-duplication, and sorting from multiple providers.
    \end{itemize}

    \subsection{Integration Testing (3rd Party Sources)}
    Targeting the \texttt{discover} package using \texttt{MockRestServiceServer}.
    \begin{itemize}
        \item \textbf{AgendaLxProviderIntegrationTest}: Verifies parsing of Lisbon's Agenda Cultural JSON.
        \item \textbf{TicketmasterProviderIntegrationTest}: Verifies parsing of Ticketmaster API responses.
        \item \textbf{SeatGeekProviderIntegrationTest}: Verifies parsing of SeatGeek API responses.
    \end{itemize}

    \subsection{REST API Integration (Controllers)}
    Comprehensive coverage using \texttt{MockMvc}.
    \begin{itemize}
        \item \textbf{AuthControllerTest}: Tests registration and login flows.
        \item \textbf{MeetingControllerTest}: Tests meeting proposal and error handling.
        \item \textbf{DiscoveryControllerTest}: Tests event discovery results and copying discovered events.
        \item \textbf{CalendarControllerTest}: Tests calendar rendering for authenticated users.
        \item \textbf{ICalControllerTest}: Verifies the public iCal feed endpoint.
    \end{itemize}

    \subsection{Security Testing}
    \begin{itemize}
        \item \textbf{SecurityTest}: Specifically verifies that public endpoints are accessible and protected endpoints redirect to login for anonymous users, while allowing access to authenticated users.
    \end{itemize}

    \subsection{Database Integration}
    \begin{itemize}
        \item \textbf{MeetingServiceIntegrationTest}: Performs a complete integration cycle from proposal to database state update.
    \end{itemize}

    \subsection{End-to-End Testing (Selenium)}
    \begin{itemize}
        \item \textbf{FullFlowE2ETest}: Performs the basic registration $\to$ login $\to$ propose journey.
        \item \textbf{ExpandedE2ETest}: Verifies discovery searching and logout flows, ensuring robustness of the browser-level interactions.
    \end{itemize}

    \section{CI/CD and Test Strategy Refactor}
    To optimize the development lifecycle, the test suite was refactored with a multi-tier execution strategy incorporating real-world CI best practices.

    \subsection{JUnit 5 Tags}
    Tests are categorized into five groups:
    \begin{itemize}
        \item \texttt{unit}: Isolated logic tests ($< 1$s).
        \item \texttt{smoke}: Critical paths required for application stability.
        \item \texttt{integration}: Tests requiring Spring context or database.
        \item \texttt{provider}: External service integration tests (mocked).
        \item \texttt{e2e}: Slow browser automation tests.
    \end{itemize}

    \subsection{Maven Configuration}
    The \texttt{pom.xml} was updated to support test filtering via properties:
    \begin{lstlisting}[language=bash]
mvn test -Dgroups="smoke"
    \end{lstlisting}

    \subsection{GitHub Actions Pipeline}
    A multi-tier pipeline was implemented in \texttt{.github/workflows/ci.yml} following industry-standard practices:

    \subsubsection{Pipeline Triggers}
    The workflow runs on three events:
    \begin{itemize}
        \item \textbf{Push to any branch}: Triggers fast feedback jobs.
        \item \textbf{Pull Request to main/master}: Triggers integration validation.
        \item \textbf{Manual dispatch} (\texttt{workflow\_dispatch}): Allows on-demand execution of the full test suite on any branch without requiring a merge.
    \end{itemize}

    \subsubsection{Concurrency Control}
    The pipeline uses GitHub's \texttt{concurrency} feature to automatically cancel outdated runs when new commits are pushed to the same branch. This prevents wasted runner minutes and ensures only the latest commit is being validated.

    \subsubsection{Job Structure and Dependencies}
    \begin{itemize}
        \item \textbf{Fast Tests} (\texttt{unit + smoke}): Runs on every push to any branch. Includes a 10-minute timeout and publishes test reports via the Dorny Test Reporter for inline GitHub UI visualization. Artifacts are uploaded for offline debugging.
        \item \textbf{Integration Tests} (\texttt{unit + smoke + integration}): Runs only on Pull Requests. Depends on fast tests passing first (\texttt{needs: fast-tests}), avoiding unnecessary compute if basic tests fail. 15-minute timeout.
        \item \textbf{Full Tests} (\texttt{all groups}): Runs only on pushes to \texttt{main} or \texttt{master}. Also depends on fast tests. 20-minute timeout. This is the production gate.
        \item \textbf{Manual Full Tests} (\texttt{all groups}): Runs exclusively via \texttt{workflow\_dispatch}. Allows developers to validate the complete suite on a feature branch before opening a PR, without polluting the branch with merge commits.
        \item \textbf{Build JAR}: Runs on every push, packaging the application and uploading the artifact. Proves the app compiles into a deployable artifact.
    \end{itemize}

    \subsubsection{Real-World Optimizations}
    \begin{itemize}
        \item \textbf{Job Dependencies}: Slow jobs (integration, full tests) only run if fast tests pass, mimicking the \textit{fail-fast} pattern used in industry pipelines.
        \item \textbf{Timeouts}: Every job has a strict timeout to prevent runaway processes from consuming GitHub Actions minutes.
        \item \textbf{Artifact Uploads}: Test reports and JAR files are preserved as downloadable artifacts, enabling post-mortem analysis without re-running the pipeline.
        \item \textbf{Test Reporter Integration}: JUnit XML results are rendered directly in the GitHub Checks UI, making failures visible without downloading logs.
    \end{itemize}

    \subsubsection{Running Tests Locally and in CI}
    \begin{lstlisting}[language=bash]
# Run all tests locally
mvn test

# Run only fast tests (feature branch feedback)
mvn test -Dgroups="smoke | unit"

# Run integration-level tests (PR validation)
mvn test -Dgroups="smoke | unit | integration"

# Run full suite (main branch or manual dispatch)
mvn test -Dgroups="smoke | unit | integration | provider | e2e"
    \end{lstlisting}

    \section{Modifications to the SUT}
    To enable effective testing, the following changes were made:
    \begin{itemize}
        \item \textbf{Discovery Providers}: Changed from manual \texttt{RestClient} instantiation to \texttt{RestClient.Builder} constructor injection. This allows Spring's \texttt{@RestClientTest} to intercept and mock HTTP calls.
        \item \textbf{Model Classes}: Added \texttt{setId(Long id)} to \texttt{User} and \texttt{Meeting}. This was necessary to mock persistence-assigned IDs in unit and REST tests.
    \end{itemize}

    \section{How to Run the Tests}
    Ensure you have Maven and a JDK (17+) installed.
    \begin{enumerate}
        \item Navigate to the \texttt{calendar} directory.
        \item Run all tests: \texttt{mvn test}
        \item To run a specific test class: \texttt{mvn test -Dtest=ClassName}
    \end{enumerate}

    \section{Conclusion}
    The resulting test suite provides high coverage across all application layers. The combination of fast unit tests and comprehensive E2E tests ensures that both business logic and user-facing features are reliable. The CI pipeline follows real-world conventions with staged execution, manual triggers, concurrency control, and artifact management, making the development workflow efficient and maintainable.

\end{document}